% Foster sweep confound illustration — for Section 5 (Discussion),
% after the first paragraph ("The experiment establishes...") and
% before the "Limitations" paragraph.
% Accompanies Figure X: two-panel plot (diversity vs n, diversity vs density).

Figure~\ref{fig:foster} illustrates the confound concretely. We ran the
same island model GA on 13 cubic symmetric graphs from the Foster
census~\cite{foster1988census} ($n = 4$ to $30$, 30 seeds each). Because
every 3-regular graph on $n$ vertices has $m = 3n/2$ edges, the cycle
rank $\beta_1 = n/2 + 1$, and the density $m/n = 3/2$---all three
quantities are affine in~$n$, hence perfectly collinear. The left panel
of Figure~\ref{fig:foster} plots diversity against~$n$ and shows an
increasing trend: ``more nodes, more diversity.'' The right panel plots
the same data against density $m/\binom{n}{2}$ and shows a decreasing
trend: ``lower density, more diversity.'' Both narratives are literally
the same regression. This is exactly the trap Theorem~\ref{thm:confound}
identifies: when $\beta_1$, $m$, and $n$ move in lockstep, any
story about cycles is equally a story about density or scale. Only
experiments at constant density---like the directed cycle experiment in
Section~\ref{sec:results}---can break the degeneracy.
